\chapter{設計方針}

本章では本研究の設計方針を述べる．3.1 節で評価全体の構成を示し，3.2 節で確信度の引き出し方式，3.3 節でデータセットの選定，3.4 節で評価指標の定義，3.5 節で実験条件を述べる．

\section{全体構成}

本研究の評価は，図 3.1 に示す流れで行う．

\begin{figure}[htbp]
\begin{center}
\begin{tabular}{c}
\fbox{\parbox{0.75\linewidth}{\centering 問題データ（問題文・正解・回答形式）}} \\
$\downarrow$ \\
\fbox{\parbox{0.75\linewidth}{\centering プロンプト生成 …… 引き出し方式 × 問題形式でテンプレートを選択}} \\
$\downarrow$ \\
\fbox{\parbox{0.75\linewidth}{\centering モデルへの問い合わせ …… 温度 0 で API を呼び出す}} \\
$\downarrow$ \\
\fbox{\parbox{0.75\linewidth}{\centering 応答の解析 …… 回答と確信度を抽出し，確信度を 0〜1 に正規化}} \\
$\downarrow$ \\
\fbox{\parbox{0.75\linewidth}{\centering 正誤判定 …… 回答形式に応じて正解と照合}} \\
$\downarrow$ \\
\fbox{\parbox{0.75\linewidth}{\centering 指標の計算 …… ECE，Brier Score，AUROC，Murphy 分解}} \\
$\downarrow$ \\
\fbox{\parbox{0.75\linewidth}{\centering 統計的検証 …… 問題の再抽出による信頼区間}} \\
\end{tabular}
\caption{評価の処理フロー}
\label{fig:3.1}
\end{center}
\end{figure}

この流れのうち，確信度の引き出し方式によって変わるのはプロンプト生成と応答の解析の 2 箇所のみである．それ以降の処理はすべての方式で共通とする．方式間の差が処理の違いに由来しないようにするためである．

\section{確信度の引き出し方式}

本研究では 3 つの方式を比較する．いずれもモデルの重みには触れず，プロンプトのみで確信度を引き出す．

\textbf{Verb.1S（回答と確信度の同時表明）} は，回答と確信度を 1 回のやりとりで同時に答えさせる方式であり，Xiong ら \cite{xiong2024llmuncertainty} が基本形として扱っている方式に相当する．プロンプトでは，回答と，その回答が正しいと思う確率を 0.0 から 1.0 の数値で答えるよう指示し，「回答:」「確信度:」の 2 行で答えるよう形式を指定する．この方式は API の呼び出しが 1 回で済むため費用と時間が小さい．一方で，回答を生成する過程と確信度を見積もる過程が同一の生成の中で混在するため，Seo ら \cite{seo2025advice}が指摘する回答非依存性の問題が生じうる．

\textbf{Verb.2S（2 段階での確信度表明）} は，回答を得たあと，そのやりとりを保持したまま別のターンで確信度を尋ねる方式であり，Tian ら \cite{tian2023justask} の 2 段階方式に相当する．第 1 ターンでは回答のみを求め，第 2 ターンでは「あなたは先ほどの問題に対して『\{第 1 ターンの回答\}』と回答しました」と自分の回答を提示したうえで，その回答が正しい確率を尋ねる．確信度の推定が回答内容に依存するよう促すことを狙った設計であり，狙いどおりの効果が得られるかは検証の対象である．

\textbf{Ling.1S（言語表現による確信度表明）} は，確信度を数値ではなく，あらかじめ定めた 5 段階の言語表現から選ばせる方式であり，Tian ら \cite{tian2023justask} の言語表現方式に相当する．各表現は表 3.1 に示す数値に対応づけて集計する．対応づける数値は，5 段階を 0 から 1 の範囲におおむね等間隔で配置しつつ，両端を 0 と 1 にしない値とした．「ほぼ確実」を確率 1，すなわち誤る可能性がないという表明として扱うのは強すぎると判断したためである．この割り当ては暫定的なものであり，値を変えた場合に指標がどの程度動くかは別途確認する必要がある（5.5 節）．

\begin{table}[htbp]
\centering
\small
\caption{言語表現と対応する数値}
\label{tab:3.1}
\begin{tabular}{|l|l|}
\hline
言語表現 & 数値 \\
\hline
ほぼ確実 & 0.95 \\
かなり自信がある & 0.80 \\
どちらともいえない & 0.50 \\
あまり自信がない & 0.25 \\
ほとんどわからない & 0.05 \\
\hline
\end{tabular}
\end{table}

3 方式は，やりとりの回数（1 回か 2 回か）と確信度の形式（数値か言語表現か）の組み合わせとして整理できる．ただし，これは 2 つの要因を厳密に統制した実験ではない．数値と言語表現では段階数や値への対応づけも同時に変わり，1 回と2 回のやりとりでは回答を生成する際の指示や文脈も変わる．また 2 回のやりとりと言語表現を組み合わせた条件を設けていないため，2 つの要因の交互作用は評価できない．本研究が比較するのは，実際に使われうる 3 つのプロンプト条件の総合的な差である．

\section{データセット}

較正の評価に用いるデータセットには 3 つの要件がある．第一に，正誤が機械的に判定できることである．判定に人手や別のモデルを介すると，判定の誤りが較正の推定に混入する．第二に，問題数が十分にあることである．較正指標は確信度を区間に分割して計算するため，各区間に一定数の標本が必要になる．第三に，日本語として自然であることである．

予備実験では日本語版 MGSM \cite{shi2023mgsm} を用いた．MGSM は算数の文章題からなるGSM8K \cite{cobbe2021gsm8k} の一部を人手で多言語に翻訳したもので，日本語版は 250 問である．解答が数値であるため正誤判定が厳密に行え，翻訳が機械ではなく人手によるため訳文の品質も一定している．一方で第三の要件は満たしていない．日本語で入力し日本語で答えさせている以上，日本語入力に対する較正の測定にはなっているが，対象は算数の文章題に限られ，日本固有の知識や文脈を含む課題には及ばない．したがって予備実験の結果を，日本語のタスク一般に広げて論じることはできない．今後は，日本語で作成されたデータセットとして JGLUE \cite{kurihara2022jglue} に含まれるJCommonsenseQA と，JMMLU による追試を検討している．

\section{評価指標}

以下，問題数を $N$，$i$ 番目の問題に対するモデルの確信度を $c_i \in [0,1]$，回答が正解なら $y_i = 1$，不正解なら $y_i = 0$ とする．

\textbf{ECE．} 確信度の値域 $[0,1]$ を $M$ 個の区間 $B_1, \dots, B_M$ に等分し，各区間に含まれる問題の集合を同じ記号で表すと，ECE は次式で定義される．

\begin{equation}
\mathrm{ECE} = \sum_{m=1}^{M} \frac{|B_m|}{N} \bigl| \mathrm{acc}(B_m) - \mathrm{conf}(B_m) \bigr|
\end{equation}

ここで $|B_m|$ は区間 $B_m$ に含まれる問題数，$\mathrm{acc}(B_m)$ は区間内の正答率，$\mathrm{conf}(B_m)$ は区間内の平均確信度である．ECE は 0 以上 1 以下の値をとり，小さいほど較正が良い．本研究では $M = 10$ とする．

\textbf{Brier Score．} 確信度と正誤の二乗誤差の平均 $\mathrm{BS} = \frac{1}{N} \sum_{i=1}^{N} (c_i - y_i)^2$ である．0 以上 1 以下の値をとり，小さいほど良い．区間分割を必要としない点で ECE と異なる．

\textbf{Murphy 分解．} Brier Score は $\mathrm{BS} = \mathrm{REL} - \mathrm{RES} + \mathrm{UNC}$ と分解できる \cite{murphy1973partition}．確信度がとりうる値を $K$ 個の区間に分け，区間 $k$ に含まれる問題数を $n_k$，区間内の平均確信度を $\bar{c}_k$，区間内の正答率を $\bar{y}_k$，全体の正答率を $\bar{y}$ とすると，各成分は次式で与えられる．

\begin{equation}
\mathrm{REL} = \frac{1}{N} \sum_{k=1}^{K} n_k (\bar{c}_k - \bar{y}_k)^2, \quad
\mathrm{RES} = \frac{1}{N} \sum_{k=1}^{K} n_k (\bar{y}_k - \bar{y})^2, \quad
\mathrm{UNC} = \bar{y}(1 - \bar{y})
\end{equation}

REL（信頼度）は各区間における確信度と正答率のずれを表し小さいほど良く，RES（識別度）は各区間の正答率が全体の正答率からどれだけ離れているかを表し大きいほど良い．UNC（不確実度）は正誤ラベルの平均のみで決まる．

この分解の扱いについて 2 点を明示しておく．第一に，上式の右辺は，各問題の確信度をその区間の平均 $\bar{c}_k$ に置き換えた予測に対する Brier Score に対応し，個々の確信度を用いた $\mathrm{BS}$ とは一般に一致しない．等号が成り立つのは確信度が区間ごとに 1 つの値しかとらない場合であり，両者の差は第 5 章で併記する．第二に，REL は各区間の正答率 $\bar{y}_k$ を含むため，UNC を切り離すことは REL を正答率から独立にすることを意味しない．区間が 1 つしかない場合，REL は $(\bar{c} - \bar{y})^2$ となり，平均確信度と平均正答率の差そのものになる．したがって本研究は Murphy 分解を，較正と正答率を切り分ける手段としてではなく，Brier Score の内訳を把握するための補助的な情報として用いる．

\textbf{AUROC．} 確信度を判別のスコアとみなし，正答を陽性，誤答を陰性としたときのROC 曲線の下側面積である．確信度の絶対的な水準ではなく順序のみを評価し，0.5 が無作為な判別に相当する．

これら 4 つの指標を併用するのは，単一の指標では結論を誤る場合があるためである．たとえば，すべての問題に同じ確信度を答えるモデルは，その値を全体の正答率に一致させれば ECE を 0 にできるが，どの問題が誤りかは全く示せない．このとき AUROC は 0.5，RES は 0 となる．

\section{実験条件}

予備実験の条件を表 3.2 に示す．

\begin{table}[htbp]
\centering
\small
\caption{予備実験の条件}
\label{tab:3.2}
\begin{tabularx}{\linewidth}{|l|X|}
\hline
項目 & 設定 \\
\hline
対象モデル & claude-sonnet-4-6 \\
温度 & 0.0 \\
データセット & 日本語版 MGSM（250 問） \\
問題形式 & 自由回答（数値） \\
引き出し方式 & Verb.1S，Verb.2S，Ling.1S \\
試行回数 & 各方式 1 回 \\
評価指標 & ECE（10 区間），Brier Score，AUROC，Murphy 分解 \\
\hline
\end{tabularx}
\end{table}

温度を 0 とするのは，同一の問題に対する出力の揺らぎを抑えるためである．出力が確率的であっても較正そのものは測定できるが，その場合は反復実験によって変動を評価する必要がある．なお温度 0 は出力が完全に同一になることを保証するものではない．対象モデルを 1 つに絞ったのは，予備実験の目的が評価基盤の妥当性の確認と方式間の差の把握にあるためである．
